\documentclass[11pt, a4paper]{article}

\usepackage[utf8]{inputenc}
\usepackage[T1]{fontenc}
\usepackage[english]{babel}
\usepackage{amsmath, amssymb, amsthm}
\usepackage{geometry}
\geometry{margin=2.5cm}
\usepackage{hyperref}

\newtheorem{theorem}{Theorem}[section]
\newtheorem{lemma}[theorem]{Lemma}
\newtheorem{conjecture}[theorem]{Conjecture}
\newtheorem{definition}[theorem]{Definition}

\title{Analysis and Partial Results on the Erd\H{o}s-Gy\'arf\'as Conjecture}
\author{Charles EDOU NZE\thanks{Charles EDOU NZE, chercheur ind\'ependant}}
\date{\today}

\begin{document}

\maketitle

\begin{abstract}
The Erd\H{o}s-Gy\'arf\'as conjecture posits that every graph with minimum degree at least 3 contains a simple cycle whose length is a power of two. In this document, we formalize the conjecture, review relevant extremal graph theory literature, and present a rigorous step-by-step structural analysis towards partial resolutions, concluding with an architecture for automated formalization in Lean 4.
\end{abstract}

\section{Introduction and Formal Statement}

Formulated in 1995 by Paul Erd\H{o}s and Andr\'as Gy\'arf\'as, the conjecture connects local degree conditions to the presence of specific cycle lengths.

\begin{definition}[Graph and Minimum Degree]
Let $G = (V, E)$ be a finite, simple, undirected graph. The minimum degree of $G$, denoted $\delta(G)$, is defined as $\min_{v \in V} \deg(v)$.
\end{definition}

\begin{definition}[Power-of-Two Cycle]
A simple cycle $C$ in $G$ is a sequence of distinct vertices $v_1, v_2, \dots, v_k$ such that $(v_i, v_{i+1}) \in E$ for $1 \le i \le k-1$, and $(v_k, v_1) \in E$. The length of $C$ is $k$. A power-of-two cycle is a cycle of length $k = 2^m$ for some integer $m \ge 2$.
\end{definition}

\begin{conjecture}[Erd\H{o}s-Gy\'arf\'as]
If a finite simple graph $G$ satisfies $\delta(G) \ge 3$, then $G$ contains a simple cycle of length $2^m$ for some integer $m \ge 2$.
\end{conjecture}

\section{Contextual Literature Search}

Extensive computational searches have been conducted. Royle demonstrated that any counterexample must have at least 17 vertices. Markstr\"om later searched for cubic counterexamples, establishing that a cubic counterexample must have at least 30 vertices, and found cubic graphs on 24 vertices whose only power-of-two cycles have length 16.

Analytically, the conjecture is known to hold for specific subclasses. Heckman and Krakovski (2013) proved it for 3-connected cubic planar graphs. Carr (2025) proved that graphs with diameter 2 and minimum degree at least 3 contain a cycle of length 4 or 8. Weaker relaxations involving average degrees have also been explored; Sudakov and Verstra\"ete (2008) showed that graphs with average degree exponential in the iterated logarithm of the number of vertices contain a cycle of length a power of two.

\section{Methodology and Structural Analysis}

We analyze the cycle space of graphs satisfying $\delta(G) \ge 3$. A standard approach in finding specific cycle lengths involves studying long paths and their closures.

\subsection{Lemma 1: Existence of a Cycle of Length at Least 4}

\begin{lemma}
Let $G = (V, E)$ be a finite simple graph with $\delta(G) \ge 3$. Then $G$ contains a simple cycle of length at least 4.
\end{lemma}

\begin{proof}
Let $P = v_1, v_2, \dots, v_k$ be a path of maximal length in $G$. Since $G$ is finite, such a maximal path exists.
Consider the vertex $v_1$. Because $P$ is maximal, all neighbors of $v_1$ must lie on the path $P$; otherwise, the path could be extended to a longer path, contradicting maximality.
Let the neighbors of $v_1$ be $v_{i_1}, v_{i_2}, \dots, v_{i_d}$, where $d = \deg(v_1) \ge \delta(G) \ge 3$.
Assume the neighbors are ordered such that $2 \le i_1 < i_2 < \dots < i_d \le k$.
Clearly, $v_{i_1} = v_2$ is adjacent to $v_1$ on the path. The edges $(v_1, v_{i_j})$ for $j \ge 2$ create cycles when combined with the subpath of $P$ from $v_1$ to $v_{i_j}$.
The length of the cycle formed by the edge $(v_1, v_{i_d})$ and the path segment from $v_1$ to $v_{i_d}$ is exactly $i_d$.
Since $d \ge 3$, the vertex $v_1$ has at least 3 distinct neighbors on the path.
The indices of these neighbors are distinct integers strictly greater than 1.
Therefore, the largest index $i_d$ must be at least $d \ge 3$. However, $v_1$ and $v_2$ account for the first two vertices.
Since the neighbors are distinct and $v_2$ is one of them, the third neighbor must have an index of at least 4.
Thus, $i_d \ge 4$.
The sequence of vertices $v_1, v_2, \dots, v_{i_d}, v_1$ forms a simple cycle of length $i_d \ge 4$.
\end{proof}

\section{Architecture for Autoformalization}

The following formalizes the definitions and the main conjecture in Lean 4.

\begin{verbatim}
import Mathlib.Combinatorics.SimpleGraph.Basic
import Mathlib.Combinatorics.SimpleGraph.Connectivity

open SimpleGraph

variable {V : Type*} [Fintype V] (G : SimpleGraph V)

def min_degree_at_least_three (G : SimpleGraph V) : Prop :=
  \forall v : V, G.degree v \ge 3

def is_power_of_two (n : \mathbb{N}) : Prop :=
  \exists m : \mathbb{N}, m \ge 2 \land n = 2^m

def has_power_of_two_cycle (G : SimpleGraph V) : Prop :=
  \exists (c : G.Walk v v), c.IsCycle \land is_power_of_two c.length

theorem erdos_gyarfas_conjecture :
  min_degree_at_least_three G \to has_power_of_two_cycle G := by
  sorry
\end{verbatim}

\end{document}
